\documentclass[10pt,letterpaper]{article}

\usepackage[
  letterpaper,
  top=0.7in,
  bottom=0.7in,
  left=0.8in,
  right=0.8in
]{geometry}
\usepackage[T1]{fontenc}
\usepackage{mathptmx}
\usepackage{amsmath}
\usepackage{amssymb}
\usepackage{graphicx}
\usepackage{booktabs}
\usepackage{array}
\usepackage{cite}
\usepackage{url}
\usepackage{float}

\setlength{\parindent}{0.14in}
\setlength{\parskip}{0pt}
\pagestyle{plain}
\renewcommand{\thesection}{\Roman{section}}
\renewcommand{\thesubsection}{\Alph{subsection}.}
\renewcommand{\thetable}{\Roman{table}}
\renewcommand{\refname}{REFERENCES}

\makeatletter
\renewcommand\section{\@startsection{section}{1}{\z@}%
  {1.3ex plus 1ex minus .2ex}%
  {0.8ex plus .2ex}%
  {\centering\normalfont\small\scshape}}
\renewcommand\subsection{\@startsection{subsection}{2}{\z@}%
  {1.0ex plus .8ex minus .2ex}%
  {0.6ex plus .2ex}%
  {\normalfont\itshape}}
\renewcommand{\fnum@table}{\textsc{Table \thetable}}
\makeatother

\begin{document}

\begin{center}
{\LARGE An investigation into Lightweight Hand Pose estimation for Edge HCI\par}
\end{center}

\vspace{0.6em}
\noindent\textbf{\textit{Abstract}}---

\vspace{0.6em}
\noindent\textbf{\textit{Keywords}}---

\section{INTRODUCTION}
what is the point of HCI?
what is the point of hand pose estimation?
what is the point of making it lightweight?



\section{LITERATURE REVIEW}.

\medskip
\noindent\textbf{Point:} The literature does not converge on a single architectural strategy for efficient pose estimation.  
\textbf{Evidence:} Existing approaches include CNN-based systems such as MediaPipe Hands~\cite{mediapipe}, lightweight attention-based models~\cite{attentionlite}, hybrid CNN-transformer designs such as EdgePose~\cite{edgepose}, temporal GCN-based tracking~\cite{handobjectgcn}, and encoder-decoder approaches such as EfficientPose~\cite{efficientpose}.  
\textbf{Explanation:} This suggests that efficient pose estimation remains an open design problem, with different model families making different tradeoffs between local feature extraction, global context, and structured joint reasoning.  
\textbf{Link:} Accordingly, it is reasonable to evaluate a DRHand-inspired architecture rather than assume that the design space has already been settled.

\medskip
\noindent\textbf{Point:} Edge deployment imposes a genuine speed-accuracy tradeoff in pose estimation.  
\textbf{Evidence:} Prior work notes that high-accuracy egocentric and 3D pose systems often rely on heavy backbones and desktop-class computation~\cite{everydayegohands}, whereas deployment-oriented methods such as MediaPipe Hands~\cite{mediapipe} and SPLite Hand~\cite{splitehand} pursue simplified architectures for constrained devices.  
\textbf{Explanation:} The literature therefore shows a persistent mismatch between models optimised for benchmark accuracy and models that are practical on embedded hardware.  
\textbf{Link:} This justifies treating runtime and deployment cost as core evaluation criteria rather than secondary implementation concerns.

\medskip
\noindent\textbf{Point:} Lightweight pose estimation methods typically achieve efficiency through deliberate simplification of the pipeline.  
\textbf{Evidence:} Reported strategies include compact backbones and lightweight attention~\cite{attentionlite}, sparse or reduced-complexity computation~\cite{splitehand}, simplified decoder heads~\cite{efficientpose}, and detector-to-tracker pipelines with cropped regions of interest~\cite{mediapipe}.  
\textbf{Explanation:} These studies imply that efficient performance usually comes from reducing unnecessary computation across multiple stages of the model rather than from a single isolated modification.  
\textbf{Link:} This provides literature support for testing whether simpler prediction targets or prediction strategies can preserve useful hand-pose performance.

\medskip
\noindent\textbf{Point:} The choice of training data is likely to affect hand-pose model generalisation.  
\textbf{Evidence:} Prior work uses both synthetic datasets such as RHD and real datasets for hand and hand-object tracking, while augmentation is often introduced to improve robustness under more realistic conditions~\cite{handobjectgcn,lightweightgcnhand,handaugment}.  
\textbf{Explanation:} These data sources differ in appearance, pose distribution, and annotation conditions, so a model trained on one domain may not generalise reliably to another. This is especially relevant when a study relies heavily on synthetic supervision but aims to operate on real camera input.  
\textbf{Link:} Therefore, dataset realism and cross-domain robustness remain important considerations in lightweight hand-pose estimation.

\medskip
\noindent\textbf{Point:} The literature leaves open whether full 21-keypoint hand representations are always necessary for lightweight edge-HCI use.  
\textbf{Evidence:} Existing lightweight hand-pose systems such as MediaPipe Hands~\cite{mediapipe}, Attention!~\cite{attentionlite}, and SPLite Hand~\cite{splitehand} focus mainly on reducing model complexity while still predicting a full hand structure.  
\textbf{Explanation:} This suggests that prior work has concentrated more on compressing the model than on questioning whether the full landmark target is required for practical HCI interaction.  
\textbf{Link:} This directly motivates the present study's comparison between reduced 10-keypoint and standard 21-keypoint prediction.


\section{TECHNICAL APPROACH}

\subsection{Project Scope}

\noindent\textbf{Point:} This study focuses on lightweight 2D RGB hand pose estimation for edge HCI using a DRHand-inspired architecture.  
\textbf{Evidence:} The experimental design in this paper centres on three evaluated components: reduced landmark prediction, dual-branch fusion, and transfer from synthetic to pseudo-labeled real-image data, as reflected in the main accuracy, branch-ablation, and transfer-learning results.  
\textbf{Explanation:} These components were chosen because they address the main practical issues encountered during the project: the need to reduce computation for embedded deployment, the need to justify the extra complexity of the dual-branch design, and the need to improve generalisation beyond synthetic-only training.  
\textbf{Link:} The scope of the study is therefore limited to questions that can be answered directly through the implemented model variants and the reported experiments.

\medskip
\noindent\textbf{Point:} The first question is whether a reduced 10-keypoint representation can retain sufficient performance for edge HCI while lowering runtime cost.  
\textbf{Evidence:} The paper includes a direct comparison between the baseline 21-keypoint model and a reduced 10-keypoint model, together with Jetson latency measurements.  
\textbf{Explanation:} This tests whether all 21 landmarks are necessary for the intended interactive setting, or whether a smaller output representation can provide a better efficiency-performance tradeoff.  
\textbf{Link:} This defines the first strand of the project scope: reduced landmark prediction for lightweight deployment.

\medskip
\noindent\textbf{Point:} The second question is whether the dual-branch fusion strategy provides a meaningful benefit over single-branch prediction.  
\textbf{Evidence:} The study reports branch-ablation results comparing fusion, heatmap-only, and coordinate-only outputs from the same trained checkpoints.  
\textbf{Explanation:} This is necessary because the DRHand-inspired design introduces extra architectural complexity, and that complexity must be justified by measurable accuracy gains rather than assumed to be beneficial.  
\textbf{Link:} This defines the second strand of the project scope: evaluating whether fusion is worth the added computation.

\medskip
\noindent\textbf{Point:} The third question is whether pseudo-labeled real-image data can improve generalisation beyond synthetic-only training.  
\textbf{Evidence:} Early experimentation relied mainly on RHD, while later experiments introduced a COCO-style hand dataset based on MediaPipe-derived labels and evaluated transfer between these two sources.  
\textbf{Explanation:} This reflects a practical constraint of the project: synthetic data was available, but access to suitable real hand-pose datasets with usable annotations was limited, and models trained only on RHD performed poorly on real camera images. The transfer-learning experiments therefore test whether imperfect pseudo-labeled real-image data can still provide a useful bridge toward more realistic deployment conditions.  
\textbf{Link:} This defines the third strand of the project scope: adaptation beyond synthetic-only supervision.

\subsection{DRHand-Inspired Model}

\noindent\textbf{Point:} The proposed model uses a shared lightweight backbone to extract features from a resized RGB hand image.  
\textbf{Evidence:} In the implementation, \texttt{HandPoseNet} is built around a single \texttt{Backbone} module that processes the input before either prediction head is applied. The backbone is composed mainly of depthwise-separable convolutions, downsampling blocks, and skip-connected upsampling stages, and outputs a compact feature representation for later prediction.  
\textbf{Explanation:} This shared backbone allows most of the computation to be reused across both prediction branches, which is important for keeping the model lightweight enough for edge-device inference.  
\textbf{Link:} Therefore, the backbone forms the efficiency-oriented core of the DRHand-inspired architecture used in this study.

Formally, the input to the model is an RGB image resized to $256 \times 256$, denoted by $x \in \mathbb{R}^{3 \times 256 \times 256}$. The shared backbone then produces a compact feature representation,
\[
F = B(x).
\]

\medskip
\noindent\textbf{Point:} The model uses two complementary prediction branches rather than a single output head.  
\textbf{Evidence:} The heatmap branch predicts one heatmap per keypoint, while the coordinate branch predicts normalized $(x,y)$ coordinates directly for each keypoint. In the implementation, the coordinate branch operates on intermediate features produced by the heatmap path rather than on a completely separate feature extractor.  
\textbf{Explanation:} This design combines the spatial localization behaviour of heatmaps with the compact direct-output behaviour of coordinate regression, allowing the model to exploit two different prediction strategies from the same backbone features.  
\textbf{Link:} As a result, the architecture can support both branch-specific evaluation and fused prediction at inference time.

From the shared representation, the heatmap branch predicts one heatmap per keypoint,
\[
\hat{H} = R_h(F), \qquad \hat{H} \in \mathbb{R}^{K \times 64 \times 64},
\]
where $K$ is the number of predicted keypoints. The coordinate branch predicts normalized landmark coordinates,
\[
\hat{P}^{coord} = R_c(F_{64}), \qquad \hat{P}^{coord} \in \mathbb{R}^{K \times 2},
\]
where $F_{64}$ denotes the intermediate feature map passed from the heatmap path into the coordinate regressor.

\medskip
\noindent\textbf{Point:} The final predicted pose can be produced by either individual branch or by a simple fusion rule.  
\textbf{Evidence:} The implementation supports heatmap-only, coordinate-only, and fusion inference modes. In fusion mode, heatmaps are decoded by argmax and then compared with the coordinate predictions on a per-joint basis, with the final choice controlled by a threshold $\alpha$ derived from the median predicted bone length.  
\textbf{Explanation:} This rule attempts to combine the strengths of the two branches while keeping post-processing simple enough for lightweight deployment.  
\textbf{Link:} Consequently, the model architecture is not limited to a single prediction strategy, which makes it suitable for the branch-ablation and fusion experiments reported later in the paper.

At inference time, the heatmap prediction for keypoint $i$ is decoded by argmax. If $(u_i^*, v_i^*)$ is the index of the peak response in $\hat{H}_i$, then the decoded heatmap coordinate is
\[
\hat{p}^{hm}_i =
\left(
\frac{u_i^*}{W-1},
\frac{v_i^*}{H-1}
\right).
\]
For fusion, the disagreement between the two predictions for joint $i$ is
\[
d_i = \left\| \hat{p}^{hm}_i - \hat{p}^{coord}_i \right\|_2.
\]
A sample-specific threshold $\alpha$ is computed as the median predicted bone length,
\[
\alpha = \mathrm{median}_{(a,b)\in \mathcal{E}}
\left\| \hat{p}^{coord}_a - \hat{p}^{coord}_b \right\|_2,
\]
where $\mathcal{E}$ is the set of hand-skeleton edges. The fused output is then defined as
\[
\hat{p}^{fused}_i =
\begin{cases}
\hat{p}^{hm}_i, & d_i < \alpha, \\
\hat{p}^{coord}_i, & \text{otherwise.}
\end{cases}
\]

\medskip
\noindent\textbf{Point:} The same architecture is used for both the standard 21-keypoint model and the reduced 10-keypoint variant.  
\textbf{Evidence:} In the implementation, the number of predicted keypoints is controlled by the output dimensionality of the heatmap and coordinate heads, allowing the same backbone-and-branch design to be instantiated for either keypoint layout.  
\textbf{Explanation:} This keeps the comparison between the 21-keypoint and 10-keypoint models fair, because the main structural difference lies in the output representation rather than in a completely different network design.  
\textbf{Link:} Therefore, the architecture provides a consistent basis for evaluating whether reducing the predicted landmark set improves the efficiency-performance tradeoff.

\medskip
\noindent\textbf{Point:} Training uses separate objectives for the two prediction branches.  
\textbf{Evidence:} The heatmap branch is supervised using Gaussian-target mean-squared error, while the coordinate branch is supervised using Wing loss.  
\textbf{Explanation:} These losses match the form of each output: the heatmap branch is trained to reproduce spatial probability maps, whereas the coordinate branch is trained to regress landmark positions directly with a loss designed for landmark localisation.  
\textbf{Link:} This training setup reflects the dual-output nature of the model and supports the later analysis of heatmap, coordinate, and fused predictions.

For the heatmap branch, each ground-truth keypoint $p_i = (x_i, y_i)$ is converted into a Gaussian target heatmap,
\[
G_i(u,v) = \exp\left(
-\frac{(u-x_i)^2 + (v-y_i)^2}{2\sigma^2}
\right),
\]
and the branch is supervised with mean-squared error,
\[
\mathcal{L}_{hm}
=
\frac{1}{KHW}
\sum_{i=1}^{K}
\sum_{u,v}
\left(\hat{H}_i(u,v) - G_i(u,v)\right)^2.
\]
The coordinate branch is supervised directly using Wing loss,
\[
\mathcal{L}_{coord}
=
\frac{1}{2K}
\sum_{i=1}^{K}
\sum_{j \in \{x,y\}}
\phi\!\left(\hat{p}^{coord}_{i,j} - p_{i,j}\right),
\]
where
\[
\phi(\delta)=
\begin{cases}
w \ln\left(1 + \frac{|\delta|}{\epsilon}\right), & |\delta| < w, \\
|\delta| - C, & \text{otherwise,}
\end{cases}
\qquad
C = w - w \ln\left(1 + \frac{w}{\epsilon}\right).
\]

\subsection{Differences to the Original DRHand Paper}



\subsection{Experimental Setup}

- dataset detail?
- implementation details?  

\subsection{Evaluation Metrics}

- how are the models evalulated?
- what is the math?
- what doe the metrics mean?

\section{RESULTS AND ANALYSIS}

\subsection{Experimental Results}
\subsection{Qualitative Results}

\subsection{Tables}

\begin{table}[H]
\centering
\footnotesize
\caption{Main quantitative comparison between the 21-keypoint baseline and the reduced 10-keypoint model on RHD. Lower SSE and EPE are better; higher PCK@0.2 is better. Values are mean $\pm$ std over five seeds.}
\label{tab:main_accuracy}
\begin{tabular}{llcccc}
\toprule
Model & Eval set & Train joints & SSE $\downarrow$ & EPE $\downarrow$ & PCK@0.2 $\uparrow$ \\
\midrule
B0 baseline & native-21 & 21 & 0.3098 $\pm$ 0.0236 & 0.0779 $\pm$ 0.0054 & 0.9112 $\pm$ 0.0071 \\
B0 baseline & shared-10 & 21 & 0.1511 $\pm$ 0.0120 & 0.0764 $\pm$ 0.0039 & 0.9101 $\pm$ 0.0071 \\
B1 reduced & native-10 & 10 & 0.1574 $\pm$ 0.0226 & 0.0706 $\pm$ 0.0054 & 0.9099 $\pm$ 0.0046 \\
\bottomrule
\end{tabular}
\end{table}

\begin{table}[H]
\centering
\footnotesize
\caption{Branch ablation for the dual-regression design. Fusion is compared against the heatmap-only and coordinate-only outputs using the same trained checkpoints. Values are mean $\pm$ std over five seeds.}
\label{tab:branch_ablation}
\begin{tabular}{llcccc}
\toprule
Model & Eval set & Mode & SSE $\downarrow$ & EPE $\downarrow$ & PCK@0.2 $\uparrow$ \\
\midrule
B0 & native-21 & fusion & 0.3098 $\pm$ 0.0236 & 0.0779 $\pm$ 0.0054 & 0.9112 $\pm$ 0.0071 \\
B0 & native-21 & heatmap & 0.3139 $\pm$ 0.0684 & 0.0989 $\pm$ 0.0335 & 0.9374 $\pm$ 0.0066 \\
B0 & native-21 & coord & 0.3134 $\pm$ 0.0229 & 0.0788 $\pm$ 0.0062 & 0.9113 $\pm$ 0.0071 \\
B1 & native-10 & fusion & 0.1574 $\pm$ 0.0226 & 0.0706 $\pm$ 0.0054 & 0.9099 $\pm$ 0.0046 \\
B1 & native-10 & heatmap & 0.2213 $\pm$ 0.0880 & 0.0991 $\pm$ 0.0288 & 0.9022 $\pm$ 0.0393 \\
B1 & native-10 & coord & 0.1599 $\pm$ 0.0218 & 0.0715 $\pm$ 0.0070 & 0.9102 $\pm$ 0.0048 \\
\bottomrule
\end{tabular}
\end{table}

\begin{table}[H]
\centering
\footnotesize
\caption{Transfer-learning results for the four training sequences evaluated on both RHD and the COCO-style hand dataset. Values are mean $\pm$ std over five seeds, and fusion-mode evaluation is shown because it is the main operating mode in the study.}
\label{tab:transfer_learning}
\resizebox{\textwidth}{!}{%
\begin{tabular}{llcccccc}
\toprule
Training sequence & Exp. & RHD SSE & RHD EPE & RHD PCK & COCO SSE & COCO EPE & COCO PCK \\
\midrule
RHD scratch & \texttt{R\_ONLY} & 0.3091 $\pm$ 0.0269 & 0.0764 $\pm$ 0.0030 & 0.9143 $\pm$ 0.0097 & 3.7488 $\pm$ 0.7985 & 0.3365 $\pm$ 0.0075 & 0.3719 $\pm$ 0.0434 \\
COCO scratch & \texttt{C\_ONLY} & 1.4194 $\pm$ 0.2619 & 0.1609 $\pm$ 0.0207 & 0.3397 $\pm$ 0.0718 & 0.0764 $\pm$ 0.0016 & 0.0612 $\pm$ 0.0049 & 0.9834 $\pm$ 0.0007 \\
COCO $\rightarrow$ RHD & \texttt{C\_TO\_R} & 0.4504 $\pm$ 0.0162 & 0.1020 $\pm$ 0.0039 & 0.8190 $\pm$ 0.0114 & 1.6426 $\pm$ 0.2293 & 0.2678 $\pm$ 0.0114 & 0.5320 $\pm$ 0.0164 \\
RHD $\rightarrow$ COCO & \texttt{R\_TO\_C} & 0.5699 $\pm$ 0.0592 & 0.1276 $\pm$ 0.0149 & 0.7690 $\pm$ 0.0221 & 0.0980 $\pm$ 0.0064 & 0.0840 $\pm$ 0.0018 & 0.9837 $\pm$ 0.0007 \\
\bottomrule
\end{tabular}
}
\end{table}

\begin{table}[H]
\centering
\footnotesize
\caption{Comparison between the MediaPipe Hand Landmarker baseline and the learned DRHand variants on RHD. The MediaPipe rows are single evaluation runs using \texttt{--hand auto}, whereas the DRHand rows are mean $\pm$ std over five seeds. Detection success is shown only for the MediaPipe runs because the learned-model evaluations do not use the same detection stage.}
\label{tab:mediapipe_baseline}
\begin{tabular}{llcccc}
\toprule
Method & Eval set & SSE & EPE & PCK@0.2 & Detect. succ. \\
\midrule
MediaPipe & native-21 & 1.7049 & 0.0833 & 0.7707 & 0.8933 \\
MediaPipe & shared-10 & 0.8599 & 0.1378 & 0.7635 & 0.8933 \\
\texttt{B0} fusion & native-21 & 0.3098 $\pm$ 0.0236 & 0.0779 $\pm$ 0.0054 & 0.9112 $\pm$ 0.0071 & n/a \\
\texttt{B0} fusion & shared-10 & 0.1511 $\pm$ 0.0120 & 0.0764 $\pm$ 0.0039 & 0.9101 $\pm$ 0.0071 & n/a \\
\texttt{B1} fusion & native-10 & 0.1574 $\pm$ 0.0226 & 0.0706 $\pm$ 0.0054 & 0.9099 $\pm$ 0.0046 & n/a \\
\bottomrule
\end{tabular}
\end{table}

\begin{table}[H]
\centering
\footnotesize
\caption{Edge-device latency summary for the main models. Forward time measures the prediction stage only, whereas end-to-end time includes image read, preprocessing, host-to-device transfer, prediction, and JSON write. Values are mean $\pm$ std over five seeds.}
\label{tab:latency}
\begin{tabular}{llccc}
\toprule
Model & Mode & Forward ms/img $\downarrow$ & End-to-end ms/img $\downarrow$ & Images/s $\uparrow$ \\
\midrule
B0 & fusion & 27.43 $\pm$ 0.45 & 36.81 $\pm$ 1.33 & 27.20 $\pm$ 0.97 \\
B0 & heatmap & 24.28 $\pm$ 0.35 & 33.19 $\pm$ 0.56 & 30.13 $\pm$ 0.50 \\
B0 & coord & 24.08 $\pm$ 0.44 & 33.17 $\pm$ 0.86 & 30.16 $\pm$ 0.77 \\
B1 & fusion & 24.59 $\pm$ 0.61 & 33.31 $\pm$ 0.71 & 30.03 $\pm$ 0.64 \\
B1 & heatmap & 24.00 $\pm$ 0.34 & 32.64 $\pm$ 0.46 & 30.64 $\pm$ 0.43 \\
B1 & coord & 23.91 $\pm$ 0.29 & 32.50 $\pm$ 0.34 & 30.77 $\pm$ 0.32 \\
\bottomrule
\end{tabular}
\end{table}

\begin{table}[H]
\centering
\footnotesize
\caption{Detailed latency breakdown for the main fused models. These metrics are available from the benchmark pipeline and help localize where runtime is spent. Values are mean $\pm$ std over five seeds.}
\label{tab:latency_breakdown}
\begin{tabular}{lcccccc}
\toprule
Model & Read ms & Prep ms & H2D ms & Forward ms & Write ms & E2E ms \\
\midrule
B0 fusion & 4.89 $\pm$ 0.64 & 2.80 $\pm$ 0.13 & 0.32 $\pm$ 0.02 & 27.43 $\pm$ 0.45 & 0.67 $\pm$ 0.09 & 36.81 $\pm$ 1.33 \\
B1 fusion & 4.49 $\pm$ 0.08 & 2.74 $\pm$ 0.05 & 0.31 $\pm$ 0.01 & 24.59 $\pm$ 0.61 & 0.52 $\pm$ 0.02 & 33.31 $\pm$ 0.71 \\
\bottomrule
\end{tabular}
\end{table}

\begin{table}[H]
\centering
\footnotesize
\caption{Fusion diagnostics for the dual-regression post-processing rule. The oracle-match rate measures how often the final fusion choice matches the lower-error branch for visible joints. Values are mean $\pm$ std over five seeds. The table includes all currently available diagnostic exports, including a shared-10 row for \texttt{B0}.}
\label{tab:fusion_diagnostics}
\begin{tabular}{llccccc}
\toprule
Model & Eval set & HM sel. & Coord sel. & $\alpha$ mean & Disagr. mean & Oracle match \\
\midrule
B0 & native-21 & 0.426 $\pm$ 0.025 & 0.574 $\pm$ 0.025 & 0.065 $\pm$ 0.007 & 0.163 $\pm$ 0.022 & 0.481 $\pm$ 0.012 \\
B0 & shared-10 & 0.407 $\pm$ 0.024 & 0.593 $\pm$ 0.024 & 0.065 $\pm$ 0.007 & 0.162 $\pm$ 0.018 & 0.474 $\pm$ 0.011 \\
B1 & native-10 & 0.501 $\pm$ 0.053 & 0.499 $\pm$ 0.053 & 0.084 $\pm$ 0.003 & 0.171 $\pm$ 0.018 & 0.529 $\pm$ 0.037 \\
\bottomrule
\end{tabular}
\end{table}

\begin{table}[H]
\centering
\footnotesize
\caption{Hyperparameter ablation relative to the baseline \texttt{B0} configuration, using native-21 fused accuracy on RHD and fused end-to-end benchmark latency. Values are mean $\pm$ std over five seeds.}
\label{tab:hyper_ablation}
\begin{tabular}{llcccc}
\toprule
Exp & Change vs. B0 & SSE $\downarrow$ & EPE $\downarrow$ & PCK@0.2 $\uparrow$ & E2E ms $\downarrow$ \\
\midrule
B0 & baseline & 0.3098 $\pm$ 0.0236 & 0.0779 $\pm$ 0.0054 & 0.9112 $\pm$ 0.0071 & 36.81 $\pm$ 1.33 \\
L1 & $\lambda_{hm}=1.5$, $\lambda_{coord}=0.5$ & 0.3583 $\pm$ 0.0534 & 0.0725 $\pm$ 0.0096 & 0.9081 $\pm$ 0.0092 & 35.77 $\pm$ 0.45 \\
L2 & $\lambda_{hm}=0.5$, $\lambda_{coord}=1.5$ & 0.3073 $\pm$ 0.0260 & 0.0812 $\pm$ 0.0022 & 0.9115 $\pm$ 0.0095 & 35.75 $\pm$ 0.49 \\
S1 & $\sigma=1.5$ & 0.3252 $\pm$ 0.0381 & 0.0707 $\pm$ 0.0055 & 0.9116 $\pm$ 0.0086 & 35.72 $\pm$ 0.25 \\
S2 & $\sigma=3.0$ & 0.3137 $\pm$ 0.0488 & 0.0833 $\pm$ 0.0058 & 0.9177 $\pm$ 0.0081 & 36.48 $\pm$ 0.53 \\
W1 & $w=5.0$ & 0.3213 $\pm$ 0.0384 & 0.0737 $\pm$ 0.0058 & 0.9150 $\pm$ 0.0099 & 36.21 $\pm$ 0.39 \\
W2 & $w=15.0$ & 0.2979 $\pm$ 0.0441 & 0.0750 $\pm$ 0.0031 & 0.9140 $\pm$ 0.0124 & 35.91 $\pm$ 0.43 \\
E1 & $\epsilon=1.0$ & 0.2994 $\pm$ 0.0591 & 0.0695 $\pm$ 0.0011 & 0.9175 $\pm$ 0.0113 & 35.84 $\pm$ 0.23 \\
E2 & $\epsilon=4.0$ & 0.3106 $\pm$ 0.0106 & 0.0873 $\pm$ 0.0056 & 0.9110 $\pm$ 0.0067 & 36.13 $\pm$ 0.52 \\
\bottomrule
\end{tabular}
\end{table}

\clearpage
\section*{QUALITATIVE FIGURES}

\begin{figure}[H]
\centering
\includegraphics[width=\textwidth]{../qual_figures/hand_frame_small_best/WIN_20260405_12_56_28_Pro_qualitative.png}
\caption{Qualitative comparison on hand\_frame\_small example 1.}
\end{figure}

\begin{figure}[H]
\centering
\includegraphics[width=\textwidth]{../qual_figures/hand_frame_small_best/WIN_20260405_12_56_33_Pro_qualitative.png}
\caption{Qualitative comparison on hand\_frame\_small example 2.}
\end{figure}

\begin{figure}[H]
\centering
\includegraphics[width=\textwidth]{../qual_figures/hand_frame_small_best/WIN_20260405_12_56_38_Pro_qualitative.png}
\caption{Qualitative comparison on hand\_frame\_small example 3.}
\end{figure}

\begin{figure}[H]
\centering
\includegraphics[width=\textwidth]{../qual_figures/hand_frame_small_best/WIN_20260405_12_56_41_Pro_qualitative.png}
\caption{Qualitative comparison on hand\_frame\_small example 4.}
\end{figure}

\begin{figure}[H]
\centering
\includegraphics[width=\textwidth]{../qual_figures/hand_frame_small_best/WIN_20260405_12_56_51_Pro_qualitative.png}
\caption{Qualitative comparison on hand\_frame\_small example 5.}
\end{figure}

\begin{figure}[H]
\centering
\includegraphics[width=\textwidth]{../qual_figures/hand_frame_small_best/WIN_20260405_12_56_57_Pro_qualitative.png}
\caption{Qualitative comparison on hand\_frame\_small example 6.}
\end{figure}

\begin{figure}[H]
\centering
\includegraphics[width=\textwidth]{../qual_figures/hand_frame_small_best/WIN_20260405_12_57_04_Pro_qualitative.png}
\caption{Qualitative comparison on hand\_frame\_small example 7.}
\end{figure}

\begin{figure}[H]
\centering
\includegraphics[width=\textwidth]{../qual_figures/hand_frame_small_best/WIN_20260405_12_57_15_Pro_qualitative.png}
\caption{Qualitative comparison on hand\_frame\_small example 8.}
\end{figure}

\begin{figure}[H]
\centering
\includegraphics[width=\textwidth]{../qual_figures/rhd_best/00143_qualitative.png}
\caption{Qualitative comparison on RHD evaluation example 00143.}
\end{figure}

\begin{figure}[H]
\centering
\includegraphics[width=\textwidth]{../qual_figures/rhd_best/00169_qualitative.png}
\caption{Qualitative comparison on RHD evaluation example 00169.}
\end{figure}

\begin{figure}[H]
\centering
\includegraphics[width=\textwidth]{../qual_figures/rhd_best/00354_qualitative.png}
\caption{Qualitative comparison on RHD evaluation example 00354.}
\end{figure}

\begin{figure}[H]
\centering
\includegraphics[width=\textwidth]{../qual_figures/rhd_best/00380_qualitative.png}
\caption{Qualitative comparison on RHD evaluation example 00380.}
\end{figure}

\begin{figure}[H]
\centering
\includegraphics[width=\textwidth]{../qual_figures/rhd_best/00477_qualitative.png}
\caption{Qualitative comparison on RHD evaluation example 00477.}
\end{figure}

\begin{figure}[H]
\centering
\includegraphics[width=\textwidth]{../qual_figures/rhd_best/00755_qualitative.png}
\caption{Qualitative comparison on RHD evaluation example 00755.}
\end{figure}

\begin{figure}[H]
\centering
\includegraphics[width=\textwidth]{../qual_figures/rhd_best/00811_qualitative.png}
\caption{Qualitative comparison on RHD evaluation example 00811.}
\end{figure}

\begin{figure}[H]
\centering
\includegraphics[width=\textwidth]{../qual_figures/rhd_best/00912_qualitative.png}
\caption{Qualitative comparison on RHD evaluation example 00912.}
\end{figure}

\clearpage
\section{CONCLUSION AND FUTURE WORK}

\clearpage
\begin{thebibliography}{00}

\bibitem{vuletic}
T. Vuletic, A. Duffy, L. Hay, C. McTeague, and G. Campbell, ``Systematic literature review of hand gestures used in human-computer interaction interfaces,'' \emph{Int. J. Human-Computer Studies}, vol. 129, pp. 47--65, 2019.

\bibitem{wachs}
J. P. Wachs, M. Kolsch, H. Stern, and Y. Edan, ``Vision-based hand-gesture applications,'' \emph{Communications of the ACM}, vol. 54, no. 2, pp. 60--71, 2011.

\bibitem{li_vr}
Y. Li, J. Huang, F. Tian, H.-A. Wang, and G.-Z. Dai, ``Gesture interaction in virtual reality,'' \emph{Virtual Reality \& Intelligent Hardware}, vol. 1, no. 1, pp. 84--112, 2019.

\bibitem{dualreg}
D. Wei, S. An, X. Zhang, J. Tian, K. A. Tsintotas, A. Gasteratos, and H. Zhu, ``Dual regression for efficient hand pose estimation,'' in \emph{Proc. IEEE Int. Conf. Robotics and Automation}, 2022.

\bibitem{mediapipe}
F. Zhang, V. Bazarevsky, A. Vakunov, A. Tkachenka, G. Sung, C.-L. Chang, and M. Grundmann, ``MediaPipe hands: On-device real-time hand tracking,'' \emph{arXiv preprint arXiv:2006.10214}, 2020.

\bibitem{srhandnet}
Y. Wang, B. Zhang, and C. Peng, ``SRHandNet: Real-time 2D hand pose estimation with simultaneous region localization,'' \emph{IEEE Trans. Image Processing}, vol. 29, pp. 2977--2986, 2019.

\bibitem{wingloss}
Z.-H. Feng, J. Kittler, M. Awais, P. Huber, and X.-J. Wu, ``Wing loss for robust facial landmark localisation with convolutional neural networks,'' in \emph{Proc. IEEE Conf. Computer Vision and Pattern Recognition}, 2018, pp. 2235--2245.

\bibitem{fasthand}
S. An, X. Zhang, D. Wei, H. Zhu, J. Yang, and K. A. Tsintotas, ``FastHand: Fast monocular hand pose estimation on embedded systems,'' \emph{Journal of Systems Architecture}, vol. 122, p. 102361, 2022.

\bibitem{attentionlite}
N. Santavas, I. Kansizoglou, L. Bampis, E. Karakasis, and A. Gasteratos, ``Attention! A lightweight 2D hand pose estimation approach,'' \emph{IEEE Sensors Journal}, vol. 21, no. 10, pp. 11488--11501, 2021.

\bibitem{handobjectgcn}
Y. Yin, C. McCarthy, and D. Rezazadegan, ``Real-time lightweight 3D hand-object pose estimation using temporal graph convolution networks,'' in \emph{AI 2024: Advances in Artificial Intelligence}, LNCS 15443, 2024, pp. 243--255.

\bibitem{everydayegohands}
A. Prakash, R. Tu, M. Chang, and S. Gupta, ``3D hand pose estimation in everyday egocentric images,'' in \emph{Proc. ECCV}, 2024.

\bibitem{splitehand}
K. H. Yeh, T. W. Hsu, and M. Sun, ``SPLite Hand: Sparsity-aware lightweight 3D hand pose estimation,'' in \emph{Proc. AICCC 2025}, Tokyo, Japan, 2025, pp. 1--10.

\bibitem{stereojetson}
E. Williams-Linera and J. M. Ramirez-Cortes, ``Stereo vision system based on the NVIDIA Jetson Nano for real-time evaluation of yoga poses,'' in \emph{Proc. IEEE}, 2024, pp. 1--7.

\bibitem{edgepose}
L. Zhang et al., ``EdgePose: Real-time human pose estimation scheme for industrial scenes,'' \emph{IEEE Access}, vol. 12, pp. 156702--156716, 2024.

\bibitem{lightweightgcnhand}
Y. Zhang et al., ``A lightweight hand attitude estimation method based on GCN feature enhancement,'' \emph{Electronics}, vol. 13, no. 10, p. 1952, 2024.

\bibitem{handaugment}
Z. Zhang, S. Xie, M. Chen, and H. Zhu, ``HandAugment: A simple data augmentation method for depth-based 3D hand pose estimation,'' \emph{arXiv preprint arXiv:2001.00702}, 2020.

\bibitem{efficientpose}
I. B. Toft, J. P. B. Nielsen, N. K. Poulsen, and A. L. Laursen, ``EfficientPose: Scalable single-person pose estimation,'' \emph{arXiv preprint arXiv:2008.09062}, 2020.

\bibitem{blazepose}
V. Bazarevsky, I. Grishchenko, K. Raveendran, T. Zhu, F. Zhang, and M. Grundmann, ``BlazePose: On-device real-time body pose tracking,'' \emph{arXiv preprint arXiv:2006.10204}, 2020.

\end{thebibliography}

\end{document}
